\section{怎样写比较简单的说明文}

说明文是把事物的形状、性质、特征、成因、关系、功用等解说清楚，或把人物的经历、特征等表述明白的一种文体。它不像记叙文那样重在用叙述和描写来写人记事，以事感人，以情动人，也不像议论文那样重在用论证方式来阐明主张，批驳谬误，以理服人；它只是通过对客观事物的介绍、说明，使读者对这一客观事物有个明晰的、完整的了解和认识，从而增长知识和技能，即是说，它只以知识教人。初学写作者应从写比较简单的说明文入手。所谓“比较简单”，是指所介绍的对象比较简单，说明的方法也比较简单，但是，如果不能掌握写说明文的要领，就会说不明白，甚至写得不象说明文。

我们该怎样来写比较简单的说明文呢？

\textbf{1. 首先对要说明的客观事物进行细致的观察和了解。}

说明事物，必须首先了解事物。了解得清楚、具体，才能表达得清楚、具体。从这个意义上说，观察是说明的前提条件，这和写记叙文要仔细观察才能写清楚、写具体的道理是一致的。不过，写说明文更要求对被说明的事物观察得仔细、准确，既要观察全面，又要抓住特征，突出重点。例如：一九八三年全国高考作文试题之一，就是要求根据一幅《这下面没有水，再换个地方挖》的漫画，写一段三百字左右的说明性文字，向没有看过此画的人介绍画的内容。实际上，这是要求写一篇比较简单的说明文。我们只要做到仔细、准确地观察画面内容，这段说明性文字是不难写好的。

这幅漫画很简单：画面的下部是水层。水层的上面是土层。土层从左往右看有五口干井——第一口有土层的四分之三深，第二、三口只有土层的一半深；第四口最深，已经接近了水层；第五口最浅，仅像个小坑。土层的右上方有一个人，身穿背心，左手拿锹，右臂上搭着件衣服，裤腿高挽，正口衔香烟，心安理得地跨步向前。画面底线外写着画题：“这下面没有水，再换个地方挖！”因此，这幅画由“水层、土层、干井、挖井人、标题”五个部分组成，重点在于“挖井人”的形态。他挖井找水，却不目标如一，坚持开掘，而浅尝辄止，半途而废，尽管下面有水，还误认为“这下面没有水”，要“再换个地方挖”，如果他还这样继续下去的话，很可能永远挖不到水。

有些考生由于观察得准确、具体、全面，也就能表达得清楚、明白、简练。然而，不少考生不是观察得不全面，漏写了水层、土层或标题；就是观察得不具体，不能对干井的数量、深浅或人的体态、神情等作具体介绍，只说“挖了几个洞，他就走了”。还有不少考生观察得不准确，不能根据画面提供的人、物、景，找出它们之间的内部联系，融会贯通此画的内容和寓意，因此把土层看成是大坝的有之，把水层看成是石油的有之，把“挖井人”看成是劳动模范或者是二流子、坏分子甚至越南特务的有之······由于观察的错误，从而导致了对画面内容和寓意的误解：有的把画面曲解成挖通大坝灌良田，有的说是地质人员到处找“石油”，或者把画面曲解成坏人挖沟搞破坏，或者说是由于吃“大锅饭”而造成不负责任的现象，或者说是越寇特工人员不让我国人民过安身日子，等等。可见，观察得片面、粗略又不准确，就不能抓住画面的特征和重点去正确理解画意，那么，写出来的说明性文字能不能给没看过画的人清楚、准确地进行分析，也就可想而知了。

看来写比较简单的说明文，关键的一步是要准确、具体、全面地观察和分析，以便了解有关事物的特征，把握它的重点和本质。这一点，务必要认真对待，仔细体会。

\textbf{2. 要说明事物的表面特征和本质特征。}

我们在认真观察和分析的基础上，就可以编写提纲，动笔写作了。具体怎么写呢？一般地说，先要说明事物的表面特征，以便让读者对该事物有一个大体的、初步的印象。

所谓事物的表面特征，是指某一事物所特有的、能与其他事物区别开来的外部标志。例如螃蟹，如果说它的表面特征是有甲壳，或者说是壳包着肉，那就没有说明白，因为这并不是螃蟹所特有的外部标志，乌龟也是有甲壳的呀，虾子也是壳包着肉的呀；还得出螃蟹有两只如钳子般的螯和四对横向排列的足，眼有柄，横着爬，才准确地说出它的表面特征，不至于和其它水生节肢动物混淆。如果是说明稍微复杂的事物，还应对该事物的总体作概括的说明，给读者一个总的轮廓，接着再对该事物的重点部分作比较详细的说明，给读者留下具体的印象。如对上面那幅画的说明，我们先对水层、土层、五口干井、挖井人和画题总说一下，然后突出介绍那个挖井人的外貌、衣着、神态和动作，这就同一般拿铁劳动的人区别开了。

说明了事物的表面特征，只能使读者对事物获得初步印象，要让读者对事物有比较深刻的认识，就要进一步说明事物的本质特征。

所谓本质特征，是指事物本身所固有的，决定事物性质、面貌和发展的重要因素。一般地说，事物的本质特征不像事物的表面特征那样容易被认识，必须仔细观察，深入分析该事物之中各部分之间的相互关系，以及该事物与他事物的关系，才能揭示明白。如上面那幅画的本质特征，必须从那个挖井人的表面特征和他与干井、土层、水层、画题的关系中才能揭示出来，即讽刺他在找水过程中所表现出的思想方法、思想作风，如浅尝辄止、半途而废、不持之以恒、不深入开掘、缺乏科学态度等。但是不少考生却将它归纳成诸如必须要具有心灵美、改革势在必行、搞破坏一定没有好下场、铁饭碗必须砸掉，等等。这就说明观察能力差，分析能力也差，就不能正确地揭示出事物的本质特征。

\textbf{3. 说明事物必须抓住中心，突出重点，围绕中心进行选材和剪裁。}

一篇说明文究竟说明什么问题，要根据被说明的对象和写作目的而定，这就形成了所要说明的中心。中心确定之后，需要围绕中心来选择材料和安排材料：凡跟中心无关的材料，一概不要；即使跟中心有关的材料，也要分清主次。凡是能深刻地说明中心的部分要详写，与中心有关但并非关键的部分要略写，不能说明中心的材料则应毫不留情地删去。例如课文《苏州园林》，全文抓住苏州园林的基本特点，即“务必使游览者无论站在哪个点上，眼前总是一幅完美的图画”，作者围绕这个中心，分别从建筑物的布局、山水的配合、花树的映衬和景致的层次四个主要方面说明苏州园林的图画美；又从每一个角落的构图、门窗的设计和雕镂、房屋的色彩等次要方面说明苏州园林的图画美。主要的部分详写，次要的部分略写，不能表现这个中心的许许多多的材料则一笔不写。这样，我们就能清楚明白，全面具体地获得关于苏州园林的知识。再如课文《看云识天气》中卷云、卷积云、积云、高积云和高层云、雨层云、积雨云的变化，与天气的变化有着密切的关系，作者就重点写，其他的云霞与天气关系不大或无关系，作者就略写或不写，这样就突出了“看云识天气”这一中心，达到让人们从云层的变化识别天气好坏的目的。又如上面那幅画的中心，是在说明那个挖井人的错误的思想方法和思想作风，因此对那个人的外形、神态和动作要着重介绍，而对五口子井、土层、水层和画题只作简略说明，至于与本画无关的别的挖井人，别的土层水层等，则一律不要写。这样就能让读者明白说明的中心，达到让人们从中汲取教训的目的。

\textbf{4. 说明一定要清楚明白，有条有理。}

叶圣陶先生说：“说明文以‘说明白了’为成功，而议论文却以‘说服他人’为成功。”要想清楚明白地对客观事物作出说明，就要按事物本身的规律性，或按人们认识事物的规律性，或把二者结合起来，有条有理地写；否则，读者就会感到杂乱无章，弄不明白你要说明的问题。

有条有理地说明事物，需要安排好说明的顺序。说明的顺序主要的有下列几种：

（1）依据构成部分的空间位置为顺序。

事物总是由各个部分组成的，从各个部分空间位置的联系中可找到一个合理的顺序。比如说，由上到下，由前到后，由外到内，由大到小，由主要到次要。有时根据实际需要，也可以采取与上述相反的顺序，把事物有条不紊地介绍清楚。如上面那幅画，你可以先从下到上地介绍画面的内容，由水层、土层、干井说到人，然后再由画面说到画外的题目；也可以先从画外的题目说到画内的内容，然后从上到下地介绍人、干井、土层和水层。这样写，就把整幅画各个部分上下内外的关系，按照一定的空间顺序自然、清楚地联系起来，显得条理分明。如果不是这样写，而是先写干井、再写水层，又写画题，接写土层，最后写人，或者先写人，再写土层，又写干井，接写画题，最后写水层，读者只会感到头绪紊乱，语无伦次，很难弄清你究竟是在说明什么意思了。

（2）以发展变化的先后时间为顺序。

事物都是在一定的时间和空间里发展变化的，事物的发展变化又是有时间先后顺序的。按照这个顺序介绍事物，可以使文章有条有理。

例如课文《人类的出现》就是按照古猿、猿人、古人、新人这四个阶段的先后顺序来写的，使读者从中看出人类的发展过程，条理非常清楚。如果要说明某种产品的制造过程，或者说明某项工艺、技术改革的程序，则应以工作程序的先后时间为顺序，依次写去，便能把握住说明的线索，写出来也清楚明白，有条不紊。例如课文《景泰蓝的制作》，按“制坯 --- 掐丝 --- 点蓝 --- 烧蓝 --- 打磨 --- 镀金”这样一个生产程序的先后时间来写，一丝不乱，层次分明。

（3）按不同类别的互相联系为顺序。

事物既有各自的特性，又有共同的属性。事物的这种特性和属性，为我们写说明文创造了划分门类的条件。只要对要说明的内容进行分析，合理归类，然后按一定的顺序加以排列，或用序数标明，或用段落加以区别，就能让读者全面地清楚地认识说明的对象。例如课文《石油的用途》，先总说石油的用途，再分说根据它的物理性质、化学性质、化学成分和其他特点人们对石油的利用，使文章条理清楚，一目了然。有的说明文要求把某个事物的表面特征、功用和本质特征都写出来，我们也可以分类说明。如按事物的形状、成因、构造、性质以及用途、使用方法，还有价值、意义等分项来写，就能说得明白。一般来说，应先说表面特征，后说功用，最后说本质特征，因为事物的功用往往是由表面特征决定的，而本质特征又是由表面特征或功用决定的。

\textbf{5. 一定要以说明的表达方式为主体，要显示语言的说明性，语言风格要质朴、平实。}

说明文并不一定要板起面孔说话，可以适当地运用叙述、描写、议论、抒情等表达方式，使文章活泼生动，甚至有点风趣，但是决不能喧宾夺主，把说明放到了次要的地位，那样就不成为说明文，而是记叙文或者议论文了。这一点必须注意。例如，要我们为《这下面没有水，再换个地方挖》这幅漫画写一段说明性文字，有些考生由于不明白说明文要以说明为主的根本特点，大写一个人挖井未成的故事，着力描写地下水如何“碧波荡漾，滔滔不绝”（？），土层又如何“黑黝黝地象一堵不可逾越的监狱的高墙”（？），挖的五口干井如何“还散发着肥沃黑土的芳香”（？）；有的考生甚至编造这个挖井人如何从日出干到太阳落坡，大写他一时得意，一时苦恼的心理活动，刻画他的种种神态，这还能算“说明性文字”吗？基本改写成一段记叙性文字了，当然不合要求。另外有些考生，则边说明边议论，而议论的篇幅超过了说明的篇幅，大挖这个挖井人的思想根源，大评他的思想方法、思想作风之不当，或者大谈怎样锲而不舍、目标如一就能成功的道理，基本写成一段议论性文字了，当然也不合要求。

写说明文，一定要显示出语言的说明性，即说得明白，让人一听就懂，一看就明，这与记叙文的语言要求生动形象、议论文的语言要求简练概括是有所不同的。说明文的说明性、客观性、知识性决定了说明文必须体现出质朴、平实的语言风格。说明文有时也用一些描写、形容等表达手法和比喻、拟人等修辞手法，但是其目的仍然是为了解说明白，而不是为了感染读者。如果不注意说明文的质朴、平实的语言风格，而过多地运用描写、形容、比喻、拟人……，则很容易弄巧成拙。我们刚才引述的某些考生在说明《这下面没有水，再换个地方挖》的画面时，着力描写和形容而得到完全相反的效果，不是很可以作为教训而引起我们的深思吗？
\newpage